\chapter{Hito 10: Construcción del Dashboard en Streamlit}

\section{Objetivo}

El objetivo de este hito es construir la primera versión del dashboard macrofinanciero utilizando Streamlit.

Hasta este momento hemos construido:

\begin{verbatim}
CSV
|
V
Python
|
V
DuckDB
|
V
staging
|
V
core
|
V
mart
\end{verbatim}

En este hito conectaremos la capa \texttt{mart} con una aplicación web interactiva.

Al finalizar este hito el estudiante será capaz de:

\begin{itemize}
\item Instalar Streamlit.
\item Conectar Streamlit con DuckDB.
\item Consultar vistas de la capa \texttt{mart}.
\item Crear indicadores KPI.
\item Crear gráficos básicos.
\item Construir una primera aplicación macrofinanciera.
\end{itemize}

\section{Resultado esperado}

Al finalizar este hito deberá existir una aplicación:

\begin{verbatim}
app/streamlit_app.py
\end{verbatim}

capaz de mostrar:

\begin{itemize}
\item Resumen macrofinanciero.
\item Evolución del tipo de cambio.
\item Indicadores bancarios.
\item Dolarización.
\item Concentración bancaria.
\end{itemize}

\section{Arquitectura del hito}

El flujo será:

\begin{verbatim}
DuckDB
|
V
mart
|
V
Streamlit
|
V
Usuario
\end{verbatim}

\section{Paso 1: Instalar Streamlit}

Con el entorno virtual activo, ejecutar:

\begin{lstlisting}[language=bash]
pip install streamlit
pip install plotly
\end{lstlisting}

Actualizar dependencias:

\begin{lstlisting}[language=bash]
pip freeze > requirements.txt
\end{lstlisting}

\section{Paso 2: Crear carpeta de aplicación}

Verificar que existe:

\begin{verbatim}
app/
\end{verbatim}

Crear el archivo:

\begin{verbatim}
app/streamlit_app.py
\end{verbatim}

\section{Paso 3: Importar librerías}

Agregar:

\begin{lstlisting}[language=Python]
import duckdb
import pandas as pd
import streamlit as st
import plotly.express as px
\end{lstlisting}

\section{Paso 4: Definir ruta de DuckDB}

Agregar:

\begin{lstlisting}[language=Python]
DB_PATH = "data/warehouse/macrofinanzas.duckdb"
\end{lstlisting}

\section{Paso 5: Crear función de consulta}

Agregar:

\begin{lstlisting}[language=Python]
def run_query(query):

```
conn = duckdb.connect(DB_PATH)

df = conn.execute(
    query
).fetchdf()

conn.close()

return df
```

\end{lstlisting}

Esta función permite consultar DuckDB desde Streamlit y devolver los resultados como un DataFrame de Pandas.

\section{Paso 6: Configurar página}

Agregar:

\begin{lstlisting}[language=Python]
st.set_page_config(
page_title="Dashboard Macrofinanciero RD",
page_icon="📊",
layout="wide"
)

st.title(
"Dashboard Macrofinanciero RD"
)

st.caption(
"Proyecto de Data Warehouse macrofinanciero con DuckDB y Streamlit"
)
\end{lstlisting}

\section{Paso 7: Cargar datos desde mart}

Agregar:

\begin{lstlisting}[language=Python]
resumen = run_query(
"""
SELECT *
FROM mart.vw_resumen_macrofinanciero
ORDER BY anio_mes
"""
)

fx = run_query(
"""
SELECT *
FROM mart.vw_tipo_cambio_diario
ORDER BY fecha
"""
)

banca = run_query(
"""
SELECT *
FROM mart.vw_indicadores_banca
ORDER BY fecha
"""
)

hhi = run_query(
"""
SELECT *
FROM mart.vw_hhi_bancario
ORDER BY anio_mes
"""
)
\end{lstlisting}

\section{Paso 8: Crear KPIs principales}

Agregar:

\begin{lstlisting}[language=Python]
ultimo = resumen.sort_values(
"anio_mes"
).tail(1)

if not ultimo.empty:

```
tasa_promedio = ultimo[
    "tasa_promedio_mes"
].iloc[0]

activos = ultimo[
    "activos_totales_sistema"
].iloc[0]

depositos = ultimo[
    "depositos_sistema"
].iloc[0]

dolarizacion = ultimo[
    "dolarizacion_depositos_sistema"
].iloc[0]

col1, col2, col3, col4 = st.columns(4)

col1.metric(
    "Tipo de cambio promedio",
    round(tasa_promedio, 2)
)

col2.metric(
    "Activos del sistema",
    f"{activos:,.0f}"
)

col3.metric(
    "Depósitos del sistema",
    f"{depositos:,.0f}"
)

col4.metric(
    "Dolarización depósitos",
    f"{dolarizacion:.2%}"
)
```

\end{lstlisting}

\section{Paso 9: Gráfico de tipo de cambio}

Agregar:

\begin{lstlisting}[language=Python]
st.subheader(
"Evolución del tipo de cambio"
)

fig_fx = px.line(
fx,
x="fecha",
y="tasa_promedio",
color="moneda_cod",
title="Tipo de cambio promedio diario"
)

st.plotly_chart(
fig_fx,
use_container_width=True
)
\end{lstlisting}

\section{Paso 10: Gráfico de indicadores bancarios}

Agregar:

\begin{lstlisting}[language=Python]
st.subheader(
"Indicadores bancarios"
)

entidades = banca[
"entidad_nombre"
].dropna().unique()

entidad_seleccionada = st.selectbox(
"Seleccione una entidad financiera",
entidades
)

banca_filtrada = banca[
banca["entidad_nombre"]
==
entidad_seleccionada
]

fig_roa = px.line(
banca_filtrada,
x="fecha",
y="roa",
title="ROA"
)

st.plotly_chart(
fig_roa,
use_container_width=True
)

fig_roe = px.line(
banca_filtrada,
x="fecha",
y="roe",
title="ROE"
)

st.plotly_chart(
fig_roe,
use_container_width=True
)
\end{lstlisting}

\section{Paso 11: Gráfico de dolarización}

Agregar:

\begin{lstlisting}[language=Python]
st.subheader(
"Dolarización bancaria"
)

fig_dolarizacion = px.line(
banca_filtrada,
x="fecha",
y=[
"dolarizacion_depositos",
"dolarizacion_cartera"
],
title="Dolarización de depósitos y cartera"
)

st.plotly_chart(
fig_dolarizacion,
use_container_width=True
)
\end{lstlisting}

\section{Paso 12: Gráfico de concentración bancaria}

Agregar:

\begin{lstlisting}[language=Python]
st.subheader(
"Concentración bancaria"
)

fig_hhi = px.line(
hhi,
x="anio_mes",
y="hhi_activos",
title="Índice HHI por activos"
)

st.plotly_chart(
fig_hhi,
use_container_width=True
)
\end{lstlisting}

\section{Paso 13: Mostrar tablas de datos}

Agregar:

\begin{lstlisting}[language=Python]
with st.expander(
"Ver resumen macrofinanciero"
):

```
st.dataframe(
    resumen,
    use_container_width=True
)
```

with st.expander(
"Ver indicadores bancarios"
):

```
st.dataframe(
    banca,
    use_container_width=True
)
```

\end{lstlisting}

\section{Paso 14: Versión completa inicial de la aplicación}

El archivo \texttt{app/streamlit_app.py} deberá contener:

\begin{lstlisting}[language=Python]
import duckdb
import pandas as pd
import streamlit as st
import plotly.express as px

DB_PATH = "data/warehouse/macrofinanzas.duckdb"

def run_query(query):

```
conn = duckdb.connect(DB_PATH)

df = conn.execute(
    query
).fetchdf()

conn.close()

return df
```

st.set_page_config(
page_title="Dashboard Macrofinanciero RD",
page_icon="📊",
layout="wide"
)

st.title(
"Dashboard Macrofinanciero RD"
)

st.caption(
"Proyecto de Data Warehouse macrofinanciero con DuckDB y Streamlit"
)

resumen = run_query(
"""
SELECT *
FROM mart.vw_resumen_macrofinanciero
ORDER BY anio_mes
"""
)

fx = run_query(
"""
SELECT *
FROM mart.vw_tipo_cambio_diario
ORDER BY fecha
"""
)

banca = run_query(
"""
SELECT *
FROM mart.vw_indicadores_banca
ORDER BY fecha
"""
)

hhi = run_query(
"""
SELECT *
FROM mart.vw_hhi_bancario
ORDER BY anio_mes
"""
)

ultimo = resumen.sort_values(
"anio_mes"
).tail(1)

if not ultimo.empty:

```
tasa_promedio = ultimo[
    "tasa_promedio_mes"
].iloc[0]

activos = ultimo[
    "activos_totales_sistema"
].iloc[0]

depositos = ultimo[
    "depositos_sistema"
].iloc[0]

dolarizacion = ultimo[
    "dolarizacion_depositos_sistema"
].iloc[0]

col1, col2, col3, col4 = st.columns(4)

col1.metric(
    "Tipo de cambio promedio",
    round(tasa_promedio, 2)
)

col2.metric(
    "Activos del sistema",
    f"{activos:,.0f}"
)

col3.metric(
    "Depósitos del sistema",
    f"{depositos:,.0f}"
)

col4.metric(
    "Dolarización depósitos",
    f"{dolarizacion:.2%}"
)
```

st.subheader(
"Evolución del tipo de cambio"
)

fig_fx = px.line(
fx,
x="fecha",
y="tasa_promedio",
color="moneda_cod",
title="Tipo de cambio promedio diario"
)

st.plotly_chart(
fig_fx,
use_container_width=True
)

st.subheader(
"Indicadores bancarios"
)

entidades = banca[
"entidad_nombre"
].dropna().unique()

entidad_seleccionada = st.selectbox(
"Seleccione una entidad financiera",
entidades
)

banca_filtrada = banca[
banca["entidad_nombre"]
==
entidad_seleccionada
]

fig_roa = px.line(
banca_filtrada,
x="fecha",
y="roa",
title="ROA"
)

st.plotly_chart(
fig_roa,
use_container_width=True
)

fig_roe = px.line(
banca_filtrada,
x="fecha",
y="roe",
title="ROE"
)

st.plotly_chart(
fig_roe,
use_container_width=True
)

st.subheader(
"Dolarización bancaria"
)

fig_dolarizacion = px.line(
banca_filtrada,
x="fecha",
y=[
"dolarizacion_depositos",
"dolarizacion_cartera"
],
title="Dolarización de depósitos y cartera"
)

st.plotly_chart(
fig_dolarizacion,
use_container_width=True
)

st.subheader(
"Concentración bancaria"
)

fig_hhi = px.line(
hhi,
x="anio_mes",
y="hhi_activos",
title="Índice HHI por activos"
)

st.plotly_chart(
fig_hhi,
use_container_width=True
)

with st.expander(
"Ver resumen macrofinanciero"
):

```
st.dataframe(
    resumen,
    use_container_width=True
)
```

with st.expander(
"Ver indicadores bancarios"
):

```
st.dataframe(
    banca,
    use_container_width=True
)
```

\end{lstlisting}

\section{Paso 15: Ejecutar Streamlit}

Desde la raíz del proyecto:

\begin{lstlisting}[language=bash]
streamlit run app/streamlit_app.py
\end{lstlisting}

Streamlit abrirá una dirección local parecida a:

\begin{verbatim}
http://localhost:8501
\end{verbatim}

\section{Paso 16: Verificar funcionamiento}

Comprobar:

\begin{itemize}
\item El dashboard abre correctamente.
\item Los KPIs aparecen en la parte superior.
\item El gráfico del tipo de cambio se visualiza.
\item El selector de entidad funciona.
\item Las tablas aparecen dentro de los expansores.
\end{itemize}

\section{Paso 17: Problemas frecuentes}

\subsection{No module named streamlit}

Instalar:

\begin{lstlisting}[language=bash]
pip install streamlit
\end{lstlisting}

\subsection{No module named plotly}

Instalar:

\begin{lstlisting}[language=bash]
pip install plotly
\end{lstlisting}

\subsection{La vista mart no existe}

Ejecutar primero:

\begin{lstlisting}[language=bash]
python etl/run_pipeline.py
\end{lstlisting}

\subsection{La base DuckDB no existe}

Verificar que existe:

\begin{verbatim}
data/warehouse/macrofinanzas.duckdb
\end{verbatim}

\section{Buenas prácticas}

Durante este hito seguiremos estas reglas:

\begin{enumerate}
\item Streamlit debe consumir la capa \texttt{mart}.
\item No consultar directamente \texttt{staging}.
\item Separar consultas SQL de visualizaciones cuando el proyecto crezca.
\item Mantener el dashboard simple en la primera versión.
\item Validar que el pipeline corre antes de abrir Streamlit.
\end{enumerate}

\section{Checklist de validación}

\begin{itemize}
\item[$\Box$] Streamlit instalado.
\item[$\Box$] Plotly instalado.
\item[$\Box$] Archivo \texttt{streamlit_app.py} creado.
\item[$\Box$] App conectada a DuckDB.
\item[$\Box$] KPIs visibles.
\item[$\Box$] Gráficos visibles.
\item[$\Box$] Filtro de entidad funcional.
\item[$\Box$] Dashboard abierto localmente.
\end{itemize}

\section{Entregable}

Al finalizar este hito el estudiante deberá disponer de:

\begin{itemize}
\item Dashboard Streamlit funcional.
\item Conexión DuckDB--Streamlit.
\item KPIs macrofinancieros.
\item Gráficos básicos.
\item Primera aplicación analítica del proyecto.
\end{itemize}

\section{Próximo hito}

En el Hito 11 prepararemos el proyecto para publicación.

Aprenderemos a:

\begin{itemize}
\item Organizar archivos para GitHub.
\item Crear un \texttt{README.md} profesional.
\item Preparar \texttt{requirements.txt}.
\item Verificar que la aplicación pueda ejecutarse en otro entorno.
\item Preparar el despliegue en Streamlit Community Cloud.
\end{itemize}
